\chapter{Python, Git, and the command line}
\label{app:tools}

\section{Create a reproducible Python environment}
From the project root:
\begin{lstlisting}[language=bash]
python3 -m venv .venv
. .venv/bin/activate
python -m pip install --upgrade pip
python -m pip install -r examples/webapp/requirements.txt
python -m unittest discover -s examples -p 'test_*.py' -v
\end{lstlisting}
On Windows PowerShell, activation uses \texttt{.venv\textbackslash Scripts\textbackslash Activate.ps1}. The environment isolates project dependencies; it does not make untrusted code safe.

\section{Run and inspect code}
Use \texttt{python path/to/file.py} to run a script. Use \texttt{python -m compileall} to check syntax without executing application behaviour. Prefer modules with functions that can be imported and tested rather than scripts whose only interface is printing.

\section{Git essentials}
\begin{lstlisting}[language=bash]
git init
git status
git add path/to/file
git commit -m "Describe the change"
git log --oneline --decorate
git diff
git switch -c feature/rescheduling
git switch main
git merge feature/rescheduling
\end{lstlisting}
A commit is a snapshot and message, not a guarantee of quality. Do not commit passwords, API keys, private data, generated virtual environments, or local databases. Use a reviewed ignore file.

\section{Debugging loop}
Reproduce the problem with the smallest input. Read the complete error message. Identify the first violated assumption rather than patching the last symptom. Add a regression test. Change one thing. Run the relevant test and the complete suite. Record the explanation in the issue or decision log.

\section{Code review questions}
\begin{enumerate}
\item Does the change satisfy a stated requirement?
\item Are inputs validated at the boundary?
\item Are permissions checked on the server?
\item What happens on duplicate, stale, missing, or malicious input?
\item Are tests meaningful and independent?
\item Does the change affect accessibility, privacy, performance, or operations?
\item Is the documentation sufficient for another student to run it?
\end{enumerate}

\section{Command-line habits}
Use a clear working directory, quote paths with spaces, inspect before deleting, and keep generated output separate from source. A shell command can have side effects; read it before running it. The build instructions in the repository are the authoritative sequence for this book.

